% =====================================================================
% The Second Moment of the Gradients:
% A Practical-Identifiability Diagnostic for Differentiable ABMs
% Two-column pre-print (ICML-style layout, self-contained).
% =====================================================================
\documentclass[10pt,twocolumn]{article}

% ---- geometry / layout ----------------------------------------------
\usepackage[letterpaper,margin=0.75in,columnsep=0.28in]{geometry}
\usepackage[T1]{fontenc}
\usepackage[utf8]{inputenc}
\usepackage{mathptmx}          % Times text + math (group house style)
\usepackage[scaled=0.9]{helvet}
\usepackage{microtype}

% ---- math -----------------------------------------------------------
\usepackage{amsmath,amssymb,amsthm,bm}
\usepackage{mathtools}

% ---- floats / graphics ----------------------------------------------
\usepackage{graphicx}
\usepackage{booktabs}
\usepackage{caption}
\captionsetup{font=small,labelfont=bf,labelsep=period}
\usepackage{enumitem}
\setlist{topsep=2pt,itemsep=1pt,parsep=0pt,leftmargin=1.4em}

% ---- color / links --------------------------------------------------
\usepackage{xcolor}
\definecolor{linkblue}{RGB}{30,72,130}
\usepackage[colorlinks=true,linkcolor=linkblue,citecolor=linkblue,
            urlcolor=linkblue,breaklinks=true]{hyperref}

% ---- bibliography ---------------------------------------------------
\usepackage[authoryear,round]{natbib}
\bibpunct{(}{)}{;}{a}{}{,}

% ---- section spacing (tighten for two-column) -----------------------
\usepackage{titlesec}
\titlespacing*{\section}{0pt}{1.4ex plus .3ex}{0.8ex}
\titlespacing*{\subsection}{0pt}{1.1ex plus .2ex}{0.5ex}
\titleformat{\section}{\normalfont\large\bfseries}{\thesection}{0.6em}{}
\titleformat{\subsection}{\normalfont\normalsize\bfseries}{\thesubsection}{0.6em}{}
\titleformat{\paragraph}[runin]{\normalfont\normalsize\bfseries}{}{0pt}{}[.]

% ---- theorem-like ---------------------------------------------------
\newtheorem{remark}{Remark}
\newtheorem{proposition}{Proposition}

% ---- macros ---------------------------------------------------------
\newcommand{\R}{\mathbb{R}}
\newcommand{\E}{\mathbb{E}}
\newcommand{\Prob}{\mathbb{P}}
\newcommand{\Fhat}{\hat{F}}
\newcommand{\Ltot}{\mathcal{L}}
\newcommand{\Xspace}{\mathcal{X}}
\newcommand{\Tspace}{\Theta}
\newcommand{\deff}{d_{\mathrm{eff}}}
\newcommand{\muref}{\mu_{\Prob_{\mathrm{ref}}}}
\newcommand{\mutheta}{\mu_{\Prob_\theta}}
% epistemic-status tag for Discussion / Future Work (visible in draft)
\newcommand{\est}[1]{\textup{\textsf{\footnotesize[#1]}}}
% verify-before-use citation flag (author-facing; remove before submission)
\newcommand{\vc}{\textcolor{red}{$^{\dagger}$}}

% wide (single-column) abstract inside a twocolumn document
\newenvironment{onecolabstract}
  {\begin{quote}\small}
  {\end{quote}}

\title{\vspace{-3ex}\bfseries The Second Moment of the Gradients:\\[0.3ex]
A Practical-Identifiability Diagnostic for\\ Differentiable Agent-Based Models}

\author{
Pietro Bicocchi\\
\normalsize Department of Computer Science, University of Oxford\\
\normalsize Institute for New Economic Thinking, University of Oxford\\
\normalsize \texttt{pietro.bicocchi@cs.ox.ac.uk}
\vspace{-2ex}
}
\date{}

\begin{document}
\twocolumn[
\maketitle
\begin{onecolabstract}
\noindent
Agent-based models are calibrated to reproduce observed data and then perturbed
to predict the effect of an intervention. This second step is sound only when the
data constrained the parameter the intervention moves, yet calibration returns a
point estimate whether it did or not, and the fit alone cannot tell the two cases
apart. Whether a parameter is constrained is the question of \emph{practical
identifiability}; answering it locally requires the curvature of the calibration
loss, which an agent-based model---a stochastic simulator with no tractable
likelihood---does not readily supply. We observe that when a differentiable
agent-based model is calibrated by minimising a maximum mean discrepancy, the
curvature is already present in the quantities the optimiser computes. Each step
evaluates, for every random seed, the gradient of that seed's loss; the mean of
these gradients drives descent, while their second moment,
$\Fhat = M^{-1}\sum_m g_m g_m^\top$, is a generalised Gauss--Newton estimate of
the loss curvature. Its eigenvectors are the parameter combinations the data
constrains; its eigenvalue spectrum ranks them. The diagnostic is therefore read
from the gradient stream the calibration emits, on the true model, at every
iterate, rather than reconstructed offline from a surrogate or from finite
differences. We track it along the calibration path on two contrasting models---the
smooth Brock--Hommes asset-pricing model and a mean-field SIR epidemic
differentiated through a discrete lockdown---and, with a falsification test built
from three non-MMD discrepancies, show that the unidentified directions it reports
belong to the simulator and not to the choice of loss. In the SIR model the
lockdown timing and strength are identifiable only in combination: to the
incidence data, an early mild lockdown is indistinguishable from a late strong
one---and that combination is precisely the counterfactual the model would be
deployed to evaluate.
\end{onecolabstract}
\vspace{1.5ex}
]

% =====================================================================
\section{Introduction}
\label{sec:intro}

Agent-based models (ABMs) describe a complex system from the bottom up, as a
population of interacting agents whose individual decisions aggregate into the
dynamics one observes: an epidemic curve, a price series, the distribution of firm
sizes \citep{farmer2009economy, bonabeau2002agent}. Their appeal for policy is
that the mechanism is explicit. Having fit an ABM to what happened, one can change
a rule an agent follows or a parameter that governs it, re-simulate, and read off
what would have happened instead---a lockdown imposed a week earlier, a higher
capital requirement, a different tax. The counterfactual is much of the reason
ABMs are built \citep{axtell2022agent}, and it rests on a step that calibration
performs but does not examine.

Calibration fixes the parameters so that the simulator reproduces the data. Because
the likelihood of an ABM is an intractable integral over every random choice every
agent makes, this is done without it: the established methods compare simulated and
observed summary statistics through approximate Bayesian computation or the
simulated method of moments \citep{grazzini2017bayesian, platt2020comparison}, or
learn the posterior directly with neural density estimators
\citep{cranmer2020frontier, dyer2024black}. What these share is that they treat the
simulator as a black box: it can be queried but not differentiated. A smaller and
more recent line makes the simulator itself differentiable, so that the calibration
loss can be minimised by gradient descent
\citep{andelfinger2023towards, chopra2023differentiable, querabofarull2025automatic}.
Rendering an ABM differentiable turns calibration from a search into an
optimisation and brings the gradient-based machinery of modern inference---variational
methods, stochastic optimisation---to bear on it \citep{querabofarull2023bayesian}.
It is the setting this paper requires, and it is one strand among several rather
than the prevailing practice.

Across all of these methods the output is a point estimate, accompanied by some
measure of fit, and the same object is returned whether or not the data determined
it. When the data constrains a parameter, the estimate is meaningful and so is any
counterfactual that moves it. When a wide range of values reproduce the data
equally well, the optimiser still stops at one of them, and nothing in the fit
records that the value was one arbitrary choice among many. A counterfactual built
on that value is then a fact about where the optimiser halted, not about the system.
The fit looks identical in the two cases; the counterfactual does not.

Whether the data constrains a parameter is the subject of \emph{practical
identifiability}, a mature theory in systems biology and econometrics
\citep{raue2009structural, rothenberg1971identification}. In its local form it is a
statement about the curvature of the calibration loss: the loss rises steeply along
directions the data constrains and stays flat along directions it does not, and the
eigendecomposition of that curvature separates the two. What the data pins down is
typically a \emph{combination} of parameters rather than any single one, and only a
direction in parameter space can express a constraint of that shape---so the
eigenvectors, not the coordinates, carry the answer \citep{gutenkunst2007universally,
transtrum2011geometry}. The obstacle for ABMs is the curvature x itself. The
classical route to it runs through the Fisher information, which requires the
likelihood the ABM lacks; the workarounds that have reached ABMs reconstruct the
curvature offline, by finite differences over re-simulations
\citep{naumann2024exploration} or by differentiating a trained surrogate in place of
the model \citep{scarrold2024computing}. Each yields a single snapshot at the fitted
point, computed by an analysis set up separately for the purpose.

We observe that in the differentiable setting the curvature need not be
reconstructed at all, because the calibration already computes it. To take one
gradient step, the differentiable pipeline simulates the model under $M$ random
seeds and evaluates, for each, the gradient of that seed's loss,
$g_m = \nabla_\theta \Ltot_m(\theta) \in \R^P$. First-order calibration is right to
use only the mean of these gradients: it is the sufficient object for choosing a
descent direction. Their second moment,
\begin{equation}
  \Fhat(\theta) \;=\; \frac{1}{M}\sum_{m=1}^{M} g_m g_m^\top \;\in\; \R^{P\times P},
  \label{eq:Fhat-intro}
\end{equation}
answers a different question---the local shape of the loss---that the optimiser had
no reason to pose. Our claim is that, for the maximum mean discrepancy (MMD) losses
standard in differentiable-ABM calibration, this second moment is a generalised
Gauss--Newton estimate of the loss curvature, and therefore that its
eigendecomposition is a practical-identifiability diagnostic. The mean summarises
where the seeds agree on a direction of descent; the second moment summarises how
they vary, and that variation is a curvature already carried in the gradient stream
the calibration emits. It is read, not reconstructed: computed on the true model,
at every iterate, at the cost of one $P\times P$ accumulation.

That the second moment has the \emph{algebraic} form of an empirical Fisher matrix
invites a specific objection, which we meet head on. \citet{kunstner2019limitations}
show that such a matrix does not in general approximate the Hessian away from a
minimum. Their argument concerns the chain that reaches curvature through the Fisher
information; ours does not use that chain. The MMD loss is the squared norm of a
residual $r(\theta) = \mutheta - \muref$, and for any loss of that form the second
moment estimates the Gauss--Newton curvature directly, without the Fisher identity
the critique relies on (Section~\ref{sec:bg:opg}). The distinction is not a
technicality: it is what makes the reading valid near the fit, exactly where the
counterfactual is evaluated.

We test the diagnostic on two models chosen to share as little as possible. The
Brock--Hommes asset-pricing model \citep{brock1998heterogeneous} is smooth and
differentiated without any surrogate, so any unidentified direction it reports must
originate in the simulator rather than in the differentiation. The mean-field SIR
model introduces the two complications a differentiable ABM meets in practice---a
discrete lockdown made differentiable by a surrogate gradient, and a transient
rather than stationary observable---and so tests whether the reading survives the
conditions under which it would actually be used. In both, a falsification test
built from three discrepancies that are \emph{not} MMD checks whether a direction
flagged as unidentified is genuinely unconstrained in the simulator or merely
invisible to the kernel. The SIR result is the one that motivates the paper: the
lockdown's timing and strength are identifiable only in combination, so that a fit
reproducing the outbreak accurately still cannot separate an early mild lockdown
from a late strong one. A per-parameter sensitivity scan of the kind used in the
current pipeline \citep{querabofarull2025automatic} certifies the model as
calibrated; the diagnostic shows that the quantity of interest is not identified.

Section~\ref{sec:background} sets the three ingredients the method needs against the
literature. Section~\ref{sec:method} constructs the diagnostic and its falsification
test. Section~\ref{sec:experiments} reports the two models, and
Sections~\ref{sec:discussion}--\ref{sec:future} discuss what an unidentified
direction means, what the reading cannot do, and what the same matrix makes
tractable next. That last point is developed in Appendix~\ref{app:preconditioning},
where the curvature the diagnostic reads is used to precondition the calibration it
diagnoses.

% =====================================================================
\section{Background}
\label{sec:background}

The diagnostic combines three ingredients: a differentiable ABM calibrated by a
distributional loss, a curvature-based notion of practical identifiability, and the
identification of the gradient second moment with that curvature. We take them in
turn. Section~\ref{sec:bg:calibration} fixes notation for the differentiable-ABM
pipeline and the MMD loss. Section~\ref{sec:bg:identifiability} states what a
curvature-based identifiability analysis measures. Section~\ref{sec:bg:sota} reviews
how curvature has been obtained for ABMs so far. Section~\ref{sec:bg:opg}
establishes the identification the diagnostic rests on.

% --------------------------------------------------------------------
\subsection{Differentiable ABMs and the MMD loss}
\label{sec:bg:calibration}

An ABM is defined by a stochastic simulator $\mathbf{f}$ that maps a parameter
$\theta \in \Tspace \subseteq \R^P$ and a random seed to an output trajectory
$x \in \Xspace$. Running the simulator draws $x \sim \Prob_\theta$ from the
distribution the parameter induces over outputs; the likelihood
$p(x \mid \theta) = \int \mathbf{f}(\theta,\xi)\,\mathrm{d}\xi$ marginalises over the
internal randomness $\xi$ and, for any model of realistic size, can be neither
written down nor evaluated \citep{platt2020comparison}. Calibration to a reference
data set therefore proceeds without the likelihood, by minimising a discrepancy
between the simulated and observed output distributions.

A \emph{differentiable} ABM \citep{querabofarull2025automatic, chopra2023differentiable}
is one implemented in an automatic-differentiation framework, so that a gradient of
aggregate outputs with respect to $\theta$ is obtained from a single backward pass.
The simulator runs unchanged in the forward pass; in the backward pass its
non-differentiable operations---branching, $\arg\max$, discrete sampling---are
replaced by smooth surrogates.\footnote{Surrogate-gradient schemes differ in how
they trade bias against variance: smoothed control flow
\citep{andelfinger2023towards}, the Gumbel-Softmax relaxation
\citep{jang2017categorical, maddison2017concrete}, the straight-through estimator
\citep{bengio2013estimating}, and the unbiased StochasticAD scheme
\citep{arya2022automatic}. Section~\ref{sec:exp:sir} reports how the diagnostic
responds to this choice.} The gradient is thus available without differencing the
black-box model.

The loss these methods minimise compares full output distributions rather than a
handful of hand-chosen summaries, and is smooth in $\theta$. The maximum mean
discrepancy (MMD) is the standard choice \citep{gretton2012kernel, dyer2024black}.
A characteristic kernel $k$ embeds each distribution $\Prob$ as a mean element
$\mu_\Prob = \E_{x\sim\Prob}[k(\cdot,x)]$ in a reproducing kernel Hilbert space, and
the loss is the squared distance between the embeddings of the simulator and the
data,
\begin{equation}
  \Ltot(\theta) \;=\; \big\|\, \mutheta - \muref \,\big\|^2
                \;=\; \|r(\theta)\|^2,
  \qquad r(\theta) = \mutheta - \muref .
  \label{eq:mmd}
\end{equation}
In the population limit this is a proper metric on distributions, zero only when the
two coincide \citep{sriperumbudur2010hilbert}; used as the loss in a
generalised-Bayes posterior it inherits robustness to misspecification
\citep{cherief2020mmd, knoblauch2022optimization}. One property of the definition is
decisive for what follows and the rest is context: the MMD is the squared norm of a
residual $r(\theta)$ that varies smoothly with $\theta$ through the simulator. A
second property bounds what the diagnostic can claim. At finite sample size and
finite kernel bandwidth, the MMD is not equally sensitive to every feature of the
output distribution, so an identifiability read from its gradients is
identifiability \emph{as the MMD sees it}; whether that agrees with identifiability
under the simulator's own dynamics is a question we test directly in
Section~\ref{sec:method:falsification}.

A single gradient step simulates the model under $M$ independent seeds and computes,
for each, the per-seed loss $\Ltot_m(\theta)$ and its gradient
$g_m = \nabla_\theta \Ltot_m(\theta) \in \R^P$. The optimiser uses these vectors
through their mean $\bar g = M^{-1}\sum_m g_m$, which points in the direction the
loss descends on average across seeds. The mean records where the seeds agree. How
far they \emph{disagree}, and along which axes, is a separate quantity---a statement
about the local shape of the loss rather than the direction of descent---which the
optimiser has no reason to form, and which
Sections~\ref{sec:bg:identifiability}--\ref{sec:bg:opg} identify with the curvature
the identifiability analysis requires.

% --------------------------------------------------------------------
\subsection{Practical identifiability as curvature}
\label{sec:bg:identifiability}

A calibrated model returns one parameter value, but that value is informative only
insofar as the data could have distinguished it from the alternatives.
Identifiability studies that distinction and divides in two. \emph{Structural}
identifiability asks whether a parameter could ever be recovered from noiseless,
unlimited data; it is a property of the model alone
\citep{cobelli1980parameter, wieland2021structural}. \emph{Practical}
identifiability asks what a finite, noisy study can actually recover
\citep{raue2009structural}. The two diverge often, and it is the practical failure
that matters for a counterfactual: a parameter the data did not constrain can be
moved, and the intervention it represents re-evaluated, with no deterioration in fit
to signal that anything is wrong.

Locally, practical identifiability is governed by how sharply the loss rises as
$\theta$ moves away from the fit---by the curvature of the loss---and the
eigendecomposition of that curvature is the natural instrument. Along an eigenvector
with a large eigenvalue the loss climbs steeply and the corresponding parameter
combination is tightly constrained; along one with a small eigenvalue the loss is
nearly flat and the combination is left almost free. That these are eigenvectors and
not individual coordinates is the essential point. A model may fix a ratio
$\theta_1/\theta_2$ tightly while leaving $\theta_1$ and $\theta_2$ each free; no
coordinate-wise analysis detects such a constraint, because only a direction in
parameter space has the right shape to express it
\citep{eisenberg2014determining, rothenberg1971identification}.

This curvature-eigenspectrum view accounts for a striking regularity. Wherever the
spectrum has been computed---across systems biology
\citep{gutenkunst2007universally}, statistical physics \citep{machta2013parameter},
and economics \citep{naumann2024exploration}---the eigenvalues span many orders of
magnitude, spaced roughly uniformly on a logarithmic scale, with no gap separating
the constrained directions from the rest. Models with this signature are called
\emph{sloppy}, and \citet{transtrum2010why, transtrum2011geometry} gave the
phenomenon its geometric reading: the outputs a model can produce form a thin
\emph{hyper-ribbon}, long along a few stiff directions and exponentially compressed
along the many sloppy ones. Two features carry forward. The identifiable object is a
small number of parameter combinations, not the parameters one at a time; and,
because the spectrum has no gap, no eigenvalue marks a natural boundary at which a
direction stops counting as identified, so any such boundary must be imposed
deliberately (Section~\ref{sec:method:bootstrap}).

Curvature is a local description, and its limits should be stated with it. It
approximates the loss basin around the fit as a quadratic and says nothing about
disconnected minima or a distant parameter setting that fits equally well; the
profile likelihood recovers that global structure, but at a cost---on the order of
$P \times G \times T$ simulator runs for $P$ parameters, a grid of $G$ points, and
an inner optimisation of $T$ steps\vc---that is prohibitive for most ABMs
\citep{raue2009structural}. Curvature is the tractable local surrogate, exact as a
description of identifiability when the basin is approximately quadratic and
degrading smoothly as it departs from that.

% --------------------------------------------------------------------
\subsection{Curvature for ABMs: the state of the art}
\label{sec:bg:sota}

Sloppiness analysis has reached ABMs only recently, and the existing work makes the
difficulty of obtaining a curvature matrix plain. \citet{naumann2024exploration}
computed curvature spectra for the Mark-0 macroeconomic ABM and several DSGE models
and confirmed that these landscapes are sloppy, with the broad, gapless spectra
found elsewhere. Lacking a likelihood to differentiate, they form the Hessian by
finite differences: each parameter is perturbed, the model re-simulated, and the
curvature estimated from the differences. The approach carries three costs for an
ABM. It requires on the order of $P^2$ re-simulations for a full
second-difference Hessian\vc; it differences \emph{stochastic} outputs, whose
seed-to-seed noise competes with the perturbation signal and must be suppressed by
averaging many runs; and it is offline, computed once at a single fitted point after
calibration has finished---a snapshot of the geometry at the optimum, detached from
the trajectory that reached it.

\citet{scarrold2024computing}\vc\ addressed the stochastic-differencing problem
directly by differentiating a \emph{surrogate} of the model rather than the model
itself: a neural network trained to imitate the simulator, whose gradients then
stand in for the true ones. This removes the differencing noise and recovers stiff
and sloppy directions cleanly, and it is the closest precedent to the present work.
The difference is where the gradients originate. A surrogate places a second,
learned model between the analysis and its object, so the geometry one measures is
the surrogate's, faithful only to the extent that the surrogate is; and fitting the
surrogate is a separate step that precedes the analysis. The gradients we use are
the model's own, produced by the differentiable simulator during calibration.

Read together, both methods reconstruct the geometry offline---by finite differences
in one case, a trained surrogate in the other---and both return a single snapshot at
the fit. Neither reads the true model at the cost of the gradients the calibration
already produces, and neither is available before convergence. Whether such a
reading is possible reduces to one question: do the per-seed gradients the pipeline
computes already contain the curvature?

% --------------------------------------------------------------------
\subsection{The second moment is a Gauss--Newton curvature}
\label{sec:bg:opg}

The disagreement among the per-seed gradients is their second moment, the
outer-product matrix of Eq.~\eqref{eq:Fhat-intro},
$\Fhat(\theta) = M^{-1}\sum_m g_m g_m^\top$. Its quadratic form along a unit vector
$v$ is $v^\top \Fhat v = M^{-1}\sum_m (v^\top g_m)^2$, the average squared derivative
of the per-seed loss along $v$: large when the seeds agree the loss moves steeply in
that direction, small when they agree it is flat. At this level $\Fhat$ describes
gradient variability. Reading identifiability from it requires the claim that, in
this setting, gradient variability \emph{is} curvature, and that claim needs care.

The matrix $\Fhat$ is an average of gradient outer products---an \emph{empirical
Fisher} in the terminology of \citet{kunstner2019limitations}, who show that such a
matrix does not in general approximate the Hessian away from a minimum. Their result
is exact and it is worth stating why it does not bite here. The empirical Fisher
approximates the Hessian only along the chain \emph{empirical Fisher}
$\to$ \emph{Fisher information} $\to$ \emph{Hessian}, whose first link holds at a
stationary point and fails away from it. An identifiability argument routed through
the Fisher information would break exactly there. Ours is not routed through it.
``Empirical Fisher'' names the algebraic form of $\Fhat$; it does not oblige us to
read $\Fhat$ through the Fisher information $F$, and the hat here is not silently
carrying that identity. We read $\Fhat$ through the generalised Gauss--Newton
structure that the residual form of the MMD supplies.

That structure is available because the MMD is a squared residual norm,
$\Ltot = \|r(\theta)\|^2$ (Eq.~\ref{eq:mmd}). For any loss of that form the Hessian
decomposes exactly as
\begin{equation}
  \nabla^2 \Ltot(\theta) \;=\; 2\,J^\top J
  \;+\; \underbrace{\textstyle\sum_i r_i\, \nabla^2 r_i}_{\to\,0 \text{ as } r\to 0},
  \qquad J = \frac{\partial r}{\partial \theta},
  \label{eq:ggn}
\end{equation}
with $J$ the Jacobian of the residual. The first term, the Gauss--Newton matrix
$2J^\top J$, is a positive-semidefinite curvature; it is exact when $r$ is linear in
$\theta$ and accurate near a good fit, where the residual and the neglected term are
both small \citep{botev2017practical, martens2020new}. Each per-seed gradient is
that seed's residual Jacobian contracted with its residual, so $\Fhat$ is a sample
estimate of $J^\top J$: the same curvature, reached without the Fisher information.
The Fisher reading, which the Kunstner critique rightly rejects, is the one that
fails away from the optimum; the Gauss--Newton reading holds near it, and it is the
one we take.

The scalar relating $\Fhat$ to $2J^\top J$ does not affect the diagnostic. Practical
identifiability is read from the eigenvectors of $\Fhat$ and from ratios of its
eigenvalues, both invariant under an overall rescaling of the matrix; the factor of
$2$, and the constant absorbed into the sample estimate, cancel. The scale does
matter when $\Fhat$ is used quantitatively rather than for its eigenstructure---for
instance as a preconditioner (Appendix~\ref{app:preconditioning}), where it must be
calibrated---and we keep it explicit there.

Two boundaries of the reading must be fixed before the diagnostic is built. The
first is that Eq.~\eqref{eq:ggn} is a near-optimum statement: the neglected term is
small when the residual is small, so $\Fhat$ estimates the loss curvature reliably
near a good fit and less so far from it. Away from the fit $\Fhat$ remains
well-defined and meaningful---it is the second moment of the gradient cloud, a
measure of local sensitivity anisotropy---but it is that, not the Hessian.
Section~\ref{sec:method} accordingly reads identifiability only where the curvature
approximation holds, and uses the earlier trajectory to confirm that the geometry
has settled rather than to read identifiability far from the fit. The second
boundary is cost: forming $\Fhat$ is an $O(MP^2)$ accumulation and its
eigendecomposition an $O(P^3)$ operation per iterate. Against the cost of simulating
the model these are negligible for the low-dimensional models we study, though the
$P^3$ term is a consideration the programme will meet as it scales to larger
parameter spaces.

% =====================================================================
\section{The Diagnostic}
\label{sec:method}

Section~\ref{sec:bg:opg} identified $\Fhat$ with the curvature of the calibration
loss. This section turns that identification into a diagnostic: how $\Fhat$ is formed
during calibration, how its eigendecomposition is read, how it is tracked along the
path with sampling uncertainty and an effective dimension, and how the reading is
falsified against discrepancies other than the MMD.

% --------------------------------------------------------------------
\subsection{Construction}
\label{sec:method:pipeline}

The diagnostic adds one step to the standard differentiable-ABM calibration loop. At
iteration $t$, with current parameter $\theta_t$:
\begin{enumerate}
  \item Sample $M$ independent seeds and simulate $M$ trajectories at $\theta_t$.
  \item Compute the per-seed loss $\Ltot_m(\theta_t)$ and, by automatic
        differentiation, the per-seed gradient
        $g_m = \nabla_\theta \Ltot_m(\theta_t)$.
  \item Form the mean $\bar g_t = M^{-1}\sum_m g_m$ and the second moment
        $\Fhat_t = M^{-1}\sum_m g_m g_m^\top$.
  \item Update $\theta_{t+1}$ with a first-order step driven by $\bar g_t$, or with
        the curvature-preconditioned step driven by $\Fhat_t$ of
        Appendix~\ref{app:preconditioning}.
\end{enumerate}
Only step~(3) is new, and it adds one $P\times P$ accumulation---the gradients
$\{g_m\}$ that step~(2) already computes are exactly its inputs. Steps~(1) and~(2)
are the existing pipeline; the diagnostic reads a quantity the pipeline emits rather
than calling the simulator again.

We take $M \ge 2P$ throughout. For $M < P$ the matrix $\Fhat_t$ is rank-deficient by
construction, its rank at most $\min(M,P)$, and its smallest eigenvalues cannot be
told from that deficiency; $M \ge 2P$ makes $\Fhat_t$ full rank with high
probability and the bootstrap of Section~\ref{sec:method:bootstrap} well-defined.
For the experiments of Section~\ref{sec:experiments}, with $P=5$, the resulting $M$
sits well within the range the differentiable-ABM literature already uses, so the
condition adds no cost.

% --------------------------------------------------------------------
\subsection{Eigendecomposition}
\label{sec:method:eigen}

Real, symmetric and positive semidefinite, $\Fhat_t$ admits the spectral
decomposition $\Fhat_t = V_t \Lambda_t V_t^\top$ with $V_t$ orthonormal and
$\Lambda_t = \mathrm{diag}(\lambda_1^{(t)}, \dots, \lambda_P^{(t)})$ ordered
$\lambda_1^{(t)} \ge \cdots \ge \lambda_P^{(t)} \ge 0$. Dropping the iterate
superscript, the $k$th eigenvalue is
\begin{equation}
  \lambda_k \;=\; v_k^\top \Fhat\, v_k
             \;=\; \frac{1}{M}\sum_{m=1}^M (v_k^\top g_m)^2 ,
\end{equation}
the average squared projection of the per-seed gradients onto $v_k$. A large
$\lambda_k$ means the gradients reliably carry a substantial component along $v_k$,
so the loss is sensitive there and the combination $v_k^\top\theta$ is well
constrained; a small $\lambda_k$ means the loss is on average flat along $v_k$ and
the combination is left essentially free. Following \citet{gutenkunst2007universally}
we call large-eigenvalue directions \emph{stiff} and small-eigenvalue directions
\emph{sloppy}.

\paragraph{A two-parameter example}
Suppose the simulator's observable depends on $\theta_1$ and $\theta_2$ only through
their sum. Then for every seed $m$ the loss $\Ltot_m$ is a function of
$\theta_1+\theta_2$ alone, and by the chain rule
\begin{equation}
  g_m \;=\; \frac{\partial \Ltot_m}{\partial(\theta_1+\theta_2)}\,(1,1)^\top
      \;\equiv\; c_m\,(1,1)^\top
\end{equation}
for a seed-dependent scalar $c_m$. Every gradient lies along $(1,1)^\top$, so
\begin{equation}
  \Fhat \;=\; \Big(\tfrac{1}{M}\textstyle\sum_m c_m^2\Big)
              \begin{pmatrix} 1 & 1 \\ 1 & 1 \end{pmatrix},
\end{equation}
whose eigenvectors are $v_1 = (1,1)/\sqrt2$ with $\lambda_1 = \tfrac{2}{M}\sum_m
c_m^2$ and $v_2 = (1,-1)/\sqrt2$ with $\lambda_2 = 0$. The diagnostic reports that
the data identifies the sum and cannot resolve the difference, without being told
that the simulator depends only on the sum: the information is in the geometry of
the gradient cloud. Scaled to $P$ dimensions and a real simulator, this is the whole
mechanism.

% --------------------------------------------------------------------
\subsection{Tracking along calibration}
\label{sec:method:trajectory}

Because $\Fhat_t$ is formed at every iterate, the diagnostic is a trajectory rather
than a single snapshot---the property that separates it from the offline analyses of
Section~\ref{sec:bg:sota}. We summarise the trajectory with one scalar and, where the
geometry moves, two structural quantities.

\paragraph{Effective dimension}
The primary summary is the number of directions the data constrains at iterate $t$,
\begin{equation}
  \deff(t) \;=\; \#\big\{k : \lambda_k^{(t)} > \tau_t\big\},
  \label{eq:deff}
\end{equation}
with the threshold $\tau_t$ fixed in Section~\ref{sec:method:bootstrap}. Of the $P$
parameters, $\deff(t)$ independent combinations are meaningfully constrained and the
remaining $P-\deff(t)$ are statistically indistinguishable from noise. Tracked across
calibration, $\deff(t)$ shows whether the identifiability structure has settled, and
is comparable across regimes, models and runs.

\paragraph{Eigenvalue and subspace evolution}
When the spectrum shifts, two finer quantities localise the change. An
order-of-magnitude jump in an individual $\lambda_k^{(t)}$ marks a region of
parameter space where the local identifiability structure changes qualitatively.
Reorientation of the constrained subspace is measured by principal angles: because
individual eigenvectors are unstable when eigenvalues are close, we track the top-$k$
subspace $\mathcal{S}_k(t) = \mathrm{span}(v_1^{(t)}, \dots, v_k^{(t)})$ and read its
rotation between iterates from the singular values of
$V_t^{(k)\top} V_{t+1}^{(k)}$ \citep{transtrum2011geometry}. Sustained large angles
indicate a subspace that is genuinely reorienting; small angles indicate stability.
The models of Section~\ref{sec:experiments} turn out to have stable geometry once
calibration settles, and there the same machinery certifies that stability rather
than tracking motion.

% --------------------------------------------------------------------
\subsection{Sampling uncertainty and the threshold}
\label{sec:method:bootstrap}

At finite $M$, $\Fhat_t$ is a stochastic estimator, and reading a small eigenvalue
as evidence of non-identifiability means separating it from sampling noise. We
quantify that noise by bootstrapping over seeds \citep{efron1979bootstrap}. At the
iterate of interest we draw $B = 500$ resamples of the seed indices with
replacement, form $\Fhat_t^{(b)} = M^{-1}\sum_{m\in\mathcal{B}_b} g_m g_m^\top$ for
each and eigendecompose; the $2.5$th and $97.5$th percentiles of
$\{\lambda_k^{(t,b)}\}_b$ give a $95\%$ interval for each $\lambda_k^{(t)}$.

The noise floor requires care. Defining it from the bootstrap interval of the
\emph{smallest} eigenvalue fails when that eigenvalue sits at the limit of
floating-point precision, as it does in the sloppiest regimes we study: the interval
then excludes zero for numerical rather than statistical reasons and would certify
directions carrying no information. We instead set the floor at a fixed multiple of
the largest eigenvalue,
\begin{equation}
  \tau_t \;=\; 10^{-8}\,\lambda_1^{(t)},
  \label{eq:tau}
\end{equation}
a relative precision floor at the float64 conditioning limit, and count a direction
as identified only if its bootstrap lower bound clears it,
\begin{equation}
  \deff(t) \;=\; \#\big\{k : \text{lower } 95\%\text{ CI of } \lambda_k^{(t)}
                 > \tau_t\big\}.
\end{equation}
The criterion is deliberately conservative: a direction must be both statistically
resolved and numerically meaningful to count, which can miss a weakly identified
direction but never invents one.

Where mid-spectrum eigenvalues are close, their eigenvectors can swap order between
resamples---an artefact of the decomposition, not instability of the geometry. We
therefore report bootstrap uncertainty for eigenvector \emph{subspaces}, as the
$97.5$th percentile of the largest principal angle between the resampled and
full-data top-$k$ subspaces. Small angles mark a reliably resolved subspace; large
angles mark one that needs larger $M$.

% --------------------------------------------------------------------
\subsection{Falsification: simulator structure or MMD artefact?}
\label{sec:method:falsification}

The eigendecomposition of $\Fhat$ diagnoses identifiability \emph{under the MMD
loss}, and on its own cannot say whether a direction it flags as sloppy is genuinely
unconstrained in the simulator or merely invisible to the kernel. The concern is
real: a diagnostic built from gradients of the MMD sees only what the MMD is
sensitive to, and a direction flat under the MMD might be constrained by some other
discrepancy, in which case the reading describes the loss and not the model.

We test this directly. At the calibrated parameter $\theta_T$, let $v_1$ and $v_P$ be
the stiffest and sloppiest eigenvectors of $\Fhat_T$. For each, and over a range of
signed step sizes $\alpha$, we perturb to $\theta_T + \alpha v$, simulate, and
compare the outputs against those at $\theta_T$ under three discrepancies, \emph{none
of them the MMD}: the differences in the first four moments; the supremum-norm
difference of the autocorrelation functions up to a model-relevant lag; and the
differences in the $1\%, 5\%, 95\%, 99\%$ quantiles. If $\Fhat$ reflects the
simulator, perturbing along $v_P$ leaves all three unchanged while perturbing along
$v_1$ moves all three, and it is the contrast between the two---not either alone---that
demonstrates it. We report curves over $\alpha$ rather than a single step, because
the rate at which outputs diverge is itself the signature of stiff versus sloppy.

The test can fail informatively. Should a perturbation along $v_P$ prove
distinguishable under one of the three discrepancies, the flagged direction is
non-identifiable under the MMD but not under the simulator's full output, and we
would report it as \emph{MMD-specific} non-identifiability rather than the stronger
simulator-intrinsic claim. Both are useful, and we report whichever the data
support.

% =====================================================================
\section{Experiments}
\label{sec:experiments}

We validate the diagnostic on two models chosen to differ in every respect that
could confound its interpretation. The Brock--Hommes asset-pricing model
(Section~\ref{sec:exp:bh}) has smooth dynamics implemented directly in an
automatic-differentiation framework, so its per-seed gradients are exact and any
non-identifiability the diagnostic reports must come from the simulator rather than
the differentiation. The mean-field SIR model (Section~\ref{sec:exp:sir}) adds the
two complications a differentiable ABM meets in practice---a discrete decision made
differentiable by a surrogate gradient, and a transient rather than stationary
observable. Both are calibrated to synthetic data generated at a known
$\theta^\star$, which supplies a ground truth: a direction the data constrains is one
along which the calibrated estimate should not be free to move, and the synthetic
setting lets us check that correspondence. The two subsections share a common
structure---trajectory, convergence eigenstructure, falsification---so that the
contrast between the models is legible directly.

% --------------------------------------------------------------------
\subsection{Brock--Hommes across the bifurcation landscape}
\label{sec:exp:bh}

The Brock--Hommes model \citep{brock1998heterogeneous} describes an asset market of
agents who switch between heterogeneous belief types according to past performance,
with the sharpness of switching set by an intensity of choice $\beta$. We adopt the
two-type specification, whose $P=5$ parameters are $\beta$ together with a trend
coefficient $g_h$ and a bias $b_h$ for each type $h \in \{1,2\}$. The model is a
canonical route to chaos: as $\beta$ increases, the price dynamics pass from a stable
steady state through periodic oscillation into deterministic chaos. We calibrate from
three starting points, one in each regime, and track the diagnostic along each
trajectory. Because the dynamics are smooth and differentiated without a surrogate,
the model isolates what the data constrains from the separate question of whether the
gradient is faithful, which we defer to Section~\ref{sec:exp:sir}.

\begin{figure*}[t]
  \centering
  \includegraphics[width=\textwidth]{images/fig2_bh_motion.png}
  \caption{\textbf{Identifiability geometry in motion (Brock--Hommes).}
  Eigenvalues $\lambda_k(t)$ of $\Fhat$ (top, with bootstrap bands) and the
  effective dimension $\deff(t)$ (bottom) across calibration, for trajectories
  initialised in the stable, periodic and chaotic regimes. The red line is the
  precision floor $\tau_t = 10^{-8}\lambda_1(t)$; the shaded region below it is
  indistinguishable from zero in float64 arithmetic and does not count toward
  $\deff$. The geometry settles within about five iterations, so the constrained
  dimension is certified throughout calibration rather than only at convergence, and
  it is regime-dependent: the stable regime constrains two directions, the periodic
  and chaotic regimes four. No bifurcation boundary is crossed in these runs.}
  \label{fig:bh_motion}
\end{figure*}

Figure~\ref{fig:bh_motion} tracks the spectrum and the effective dimension across
calibration. The geometry settles within roughly five iterations and is thereafter
stationary, so the diagnostic does not merely describe the converged point but
certifies at every iterate how many parameter combinations the data constrains. We
had allowed that a trajectory crossing a bifurcation boundary might leave a signature
in the spectrum (Section~\ref{sec:method:trajectory}); in these runs the optimiser
stays within the regime of its initialisation, no boundary is crossed, and we report
no crossing signature rather than construct one. What the trajectory does establish
is the more basic fact that $\deff$ is stable and certifiable throughout the fit.

The constrained dimension is itself regime-dependent, and by more than a constant. By
the criterion of Eq.~\eqref{eq:deff}, the stable regime resolves two directions while
the periodic and chaotic regimes resolve four, the chaotic spectrum's fifth
eigenvalue lying close enough to the floor that $\deff$ intermittently admits it.
This has a dynamical reading, which we offer as such. Below the first bifurcation the
long-run behaviour is a fixed point the price reaches almost regardless of the fine
parameter values, so the data carries little information about them and most of the
spectrum collapses; as the dynamics acquire periodic and then chaotic structure, a
larger share of parameter space leaves a measurable imprint on the trajectory and
more directions become constrained. The diagnostic thus recovers, along the
calibration path, a regime-dependent refinement of the single-point sloppiness
reported offline by \citet{naumann2024exploration}: not that the model is uniformly
sloppy, but that how much of it the data can see depends on the dynamical regime the
calibration sits in.

\begin{figure}[t]
  \centering
  \includegraphics[width=\linewidth]{images/fig3_eigenvectors.png}
  \caption{\textbf{Eigenvector structure at convergence (Brock--Hommes).}
  $|V_{jk}|$ for each regime, parameters against eigenvectors ordered stiff to
  sloppy; the greyscale strip shows $\log_{10}\lambda_k$ (dark~$=$~stiff). The
  stiffest direction is in every regime the symmetric bias combination
  $(b_1{+}b_2)/\sqrt2$; the symmetric trend combination is next. The between-type
  difference combinations and the intensity of choice $\beta$ occupy the sloppy tail,
  with $\beta$ the sloppiest direction outright in the periodic and chaotic regimes.
  In the stable regime the two sloppiest eigenvectors are near-degenerate and are
  read as a subspace.}
  \label{fig:bh_eigenvectors}
\end{figure}

What the data does and does not constrain is read from the eigenvectors
(Figure~\ref{fig:bh_eigenvectors}). One structure recurs across all three regimes:
the stiffest direction is the symmetric bias combination $(b_1{+}b_2)/\sqrt2$, the
symmetric trend combination is next, and the sloppy tail holds the between-type
\emph{difference} combinations together with $\beta$---the sloppiest direction
outright in the periodic and chaotic regimes. The reading admits a clean economic
interpretation, which we give as interpretation and not as a derived fact: MMD
calibration pins down the market-average belief structure far more tightly than
either the heterogeneity \emph{between} belief types or the rate at which agents
switch between them. This asymmetry is invisible to a per-parameter analysis, which
would see five separately weak or strong coordinates; only the eigendecomposition
exposes that the constrained objects are particular combinations. One caveat sharpens
it: in the stable regime the two sloppiest eigenvectors are near-degenerate and
exchange order under resampling, so they are properly a two-dimensional sloppy
subspace rather than two directions---the case the subspace tracking of
Section~\ref{sec:method:bootstrap} is built for.

\begin{figure}[t]
  \centering
  \includegraphics[width=\linewidth]{images/fig4_falsification.png}
  \caption{\textbf{Falsification (Brock--Hommes).} Output discrepancy against signed
  perturbation $\alpha$ along the stiffest (top) and sloppiest (bottom) eigenvector
  of $\Fhat_T$, under three discrepancies none of which is the MMD. The stiff
  direction moves all three steeply and roughly symmetrically; the sloppy direction
  moves none of them, the tail quantiles included. The flagged non-identifiability is
  therefore a property of the simulator, not a blind spot of the MMD loss.}
  \label{fig:bh_falsification}
\end{figure}

A diagnostic assembled from gradients of the MMD could in principle report only the
blind spots of the MMD. The falsification test of
Section~\ref{sec:method:falsification} settles this, and
Figure~\ref{fig:bh_falsification} reports the result. Perturbing the calibrated
parameter along the stiffest eigenvector raises all three non-MMD discrepancies
steeply and roughly symmetrically; perturbing along the sloppiest leaves all three
essentially flat, the tail quantiles included---the discrepancy most likely to
register structure the kernel underweights. The sloppy direction is thus invisible
not only to the MMD but to three discrepancies that never saw the kernel, and it is
the contrast between the two rows that licenses the conclusion: for this model the
diagnostic measures intrinsic structure of the simulator, not its own blind spots.
Whether the reading survives truncation of the gradient horizon $H$ is taken up in
Appendix~\ref{app:technical} (Figure~\ref{fig:horizon}): the Brock--Hommes spectrum
grows in magnitude with $H$ but keeps its hierarchy, so the qualitative reading is
robust to truncation. That robustness does not survive the move to a transient model.

% --------------------------------------------------------------------
\subsection{Mean-field SIR under surrogate-gradient bias}
\label{sec:exp:sir}

The mean-field SIR model places the diagnostic where a differentiable ABM is
genuinely difficult. The simulator is a stochastic SIR epidemic with $P=5$
parameters: the transmission rate $\beta$, the recovery rate $\gamma$, the initial
infected fraction $I_0$, and a lockdown parameterised by its timing $t_{\mathrm{lock}}$
and strength $f_{\mathrm{lock}}$, which multiplies the transmission term over a window.
The lockdown is a discrete intervention on the transmission dynamics, made
differentiable by a Gumbel-Sigmoid surrogate in the tangent pass; the model therefore
exercises both surrogate-gradient bias and the truncated gradient horizon a transient
epidemic forces. We calibrate to a synthetic outbreak with basic reproduction number
$R_0 \approx 2.4$, peaking near day~17, in which the lockdown lowers peak incidence by
roughly $57\%$---an intervention large enough that its effect is exactly the kind of
counterfactual a calibrated model would be used to assess.

\begin{figure*}[t]
  \centering
  \includegraphics[width=\textwidth]{images/fig5_sota_contrast.png}
  \caption{\textbf{Per-parameter scan versus eigenstructure (mean-field SIR).}
  Left: per-parameter first-order sensitivity $\|\partial x/\partial\theta_k\|$, in
  the manner of the individual-parameter scan of \citet{querabofarull2025automatic}.
  Right: $|V_{jk}|$ of $\Fhat$ on the same parameter axis, eigenvectors ordered stiff
  to sloppy. The epidemic parameters $\beta,\gamma,I_0$ each align with a single
  eigenvector and are individually identified; the lockdown timing $t_{\mathrm{lock}}$
  and strength $f_{\mathrm{lock}}$ appear only as an entangled pair in the two
  sloppiest eigenvectors (boxed). The scan reports the lockdown parameters as
  individually weak; the eigenstructure reports their \emph{combination} as the
  least-identified object in the model. That combination is the lockdown
  counterfactual.}
  \label{fig:sir_contrast}
\end{figure*}

The central result is the contrast in Figure~\ref{fig:sir_contrast}, which sets a
per-parameter sensitivity scan of the kind used in the current differentiable-ABM
pipeline \citep{querabofarull2025automatic} beside the eigenstructure of $\Fhat$. The
epidemic parameters $\beta$, $\gamma$ and $I_0$ each align almost exactly with their
own eigenvector and are individually identified. The lockdown timing and strength
never appear in isolation: they are entangled across the two sloppiest eigenvectors,
$v_4$ and $v_5$, so that only their combination is constrained, with condition number
$\lambda_1/\lambda_P \sim 10^{5.9}$. The two analyses read the same parameters
oppositely. The per-parameter scan reports $t_{\mathrm{lock}}$ and $f_{\mathrm{lock}}$
as individually weak but present; the eigenstructure reports their combination as the
least-identified object in the model. The equifinal reading is concrete: to the
incidence data, a lockdown imposed later but enforced more strongly is
indistinguishable from one imposed earlier but more weakly. This is where the
introduction's warning lands---the lockdown effect, the counterfactual the model
exists to evaluate, lies inside the subspace the incidence data does not constrain. A
per-parameter scan would certify the model as adequately calibrated; the diagnostic
shows that the very quantity of interest is unidentified.

That this non-identifiability belongs to the epidemic and not the kernel is confirmed
by the falsification test (Appendix~\ref{app:technical},
Figure~\ref{fig:sir_falsification}): the sloppy lockdown combination is suppressed by
twenty to thirty times relative to the stiff direction under all three non-MMD
discrepancies. The stiff direction's response there is markedly asymmetric, and the
asymmetry is genuine rather than an artefact---driving the parameter toward
suppression removes the outbreak entirely, after which the discrepancy can grow no
further and saturates, since an epidemic cannot be more than fully suppressed.

Whether the biased gradient distorts the diagnostic separates cleanly into the two
sources of bias, both treated in Appendix~\ref{app:technical}, and the result runs
against the naive expectation. The surrogate is the more benign source: the
stiff/sloppy partition and the sloppy subspace are robust to the choice of surrogate
(Gumbel-Sigmoid versus straight-through, Figure~\ref{fig:surrogate}, agreeing to
$3$--$6^\circ$ on the sloppy directions), while the orientation of the stiff
directions rotates by about $28^\circ$ and the absolute eigenvalue magnitudes shift
by roughly two orders. The surrogate bias thus contaminates the best-identified
directions rather than the worst---a sharper caution than the limitation of
Section~\ref{sec:discussion}, because it touches the directions one would most want to
trust. The horizon is the decisive source (Figure~\ref{fig:horizon}). Because the
epidemic is an early transient and the truncated gradient retains only the final
steps, every horizon short of the full trajectory sees only the post-epidemic tail
and the spectrum collapses to the floor. The SIR eigenstructure of
Figure~\ref{fig:sir_contrast} is meaningful only at full horizon, and the diagnostic
is fragile to a truncation practitioners routinely apply---a fragility the smooth,
persistent Brock--Hommes dynamics did not show.

Across both models the diagnostic does the same three things. It surfaces equifinal
parameter combinations that no per-parameter scan can represent; it tracks them along
the path and certifies an effective dimension throughout calibration; and, under a
falsification test it could have failed, it reports non-identifiabilities that are
properties of the simulator rather than of the loss. What it does not establish---and
what the lockdown result makes pressing---is that a direction the data \emph{can}
constrain corresponds to anything real in the world. We turn to that limit next.

% =====================================================================
\section{Discussion}
\label{sec:discussion}

The diagnostic answers one question---which parameter combinations the calibration
data constrains---and its value and its limits both follow from the precise form of
that question.

\paragraph{What an unidentified direction means}
A sloppy eigenvector is a combination of parameters the loss cannot separate, and its
practical import is that any counterfactual moving $\theta$ along it changes a stated
quantity while leaving the fit intact \est{established}. The SIR lockdown is the
sharp case: because $t_{\mathrm{lock}}$ and $f_{\mathrm{lock}}$ enter the sloppy
subspace only in combination, a policy question posed in terms of one of them
alone---``what if the lockdown had come a week earlier?''---has no answer the
incidence data can support, even though the model fits that data well and returns a
definite estimate for each parameter. The diagnostic does not repair this; it makes
it visible before the counterfactual is trusted. This is the floor it sets: a
claim about a parameter the data cannot identify is not a claim the calibration
licenses, whatever the quality of the fit.

\paragraph{Identifiability as the MMD sees it}
The eigenstructure is read from gradients of the MMD, so what it reports is
identifiability under that loss \est{established}. The falsification test is what
keeps this from being a limitation in disguise: by perturbing along the flagged
directions and measuring three discrepancies the kernel never saw, it distinguishes a
direction the simulator leaves unconstrained from one the MMD merely underweights.
In both our models the sloppy directions passed that test, so the reading is
simulator-intrinsic rather than kernel-specific. We do not claim this holds for every
model and kernel \est{open question}; we claim the test decides it case by case, and
that a diagnostic which can be falsified and is not is worth more than one that cannot
be.

\paragraph{Locality}
The curvature reading is a local, near-optimum statement (Section~\ref{sec:bg:opg}).
It describes the loss basin around the fit and is silent about a distant parameter
setting that fits equally well, or a disconnected minimum \est{established}. The
trajectory view mitigates but does not remove this: tracking $\Fhat_t$ from
initialisation shows whether the geometry has settled, which guards against reading a
transient as structure, but a settled local geometry is still local. Global
identifiability remains the province of the profile likelihood, at a cost the
diagnostic is designed to avoid.

\paragraph{The bias inverts the usual caution}
The surrogate-gradient result is the one we would most flag to a practitioner. The
sloppy subspace---the part of the answer a user is most likely to act on, since it
tells them what \emph{not} to trust---is robust to the surrogate; the stiff
directions and the eigenvalue magnitudes are not \est{established, on two surrogates}.
The bias therefore sits in the best-identified directions. For the qualitative
question the diagnostic is built for---which combinations are unconstrained---this is
reassuring. For any quantitative use of the stiff eigenvalues, including
preconditioning, it is a caution that must be carried explicitly. Whether it
generalises beyond the two surrogates we tried is open \est{open question}.

\paragraph{Relation to prior ABM sloppiness analyses}
The finding that ABM loss landscapes are sloppy is not ours; it is
\citet{naumann2024exploration} and, via surrogates, \citet{scarrold2024computing}.
What the emission reading adds is that the curvature need not be reconstructed offline
in the differentiable setting: it is present in the gradient stream calibration
already produces, and so can be read on the true model and at every iterate rather
than once at the end on a surrogate or by finite differences \est{established, this
work}. The regime-dependence of $\deff$ in Brock--Hommes
(Figure~\ref{fig:bh_motion}) is a refinement the along-the-path view makes visible and
a single-point analysis cannot.

% =====================================================================
\section{Future Work}
\label{sec:future}

Once the reading is validated on tractable benchmarks, the same matrix becomes an
instrument that can be pointed at harder questions.

\paragraph{From diagnosis to data design}
The sloppy eigenvectors name the combinations the data cannot constrain, which
immediately raises the constructive question: what additional summary statistic, or
what micro-level observation, would add rank to $\Fhat$ along those directions? This
reframes identifiability as a design problem---not only diagnosing what is
unidentified but prescribing what would identify it \est{plausible; the machinery
supports it, the experiment is not yet done}. For the SIR lockdown, the concrete form
is whether an observable other than aggregate incidence---a second wave, an
age-stratified curve---separates $t_{\mathrm{lock}}$ from $f_{\mathrm{lock}}$.

\paragraph{Structure and calibration geometry}
Loosely in the spirit of the loss-landscape studies of \citet{li2018visualizing} for
neural networks, one can ask how structural properties of an ABM---population size
$N$, agent heterogeneity, interaction topology---systematically shape the $\Fhat$
spectrum \est{open question}. As $N$ grows, the law of large numbers damps
idiosyncratic variation and the aggregate output should grow less sensitive to
individual perturbations, contracting the spectrum; richer heterogeneity might add
stiff directions by enriching the aggregate dynamics, or sloppy ones by introducing
new equifinality between type configurations \est{speculation}. Whether these effects
are systematic and predictable is empirical, and the emission reading makes it cheap
to ask.

\paragraph{Beyond identifiability}
Identifiability is a floor, not a guarantee. A well-identified but misspecified model
produces confident estimates of parameters that correspond to nothing in the world
\est{established, general}. The diagnostic certifies that the data constrained a
combination; it says nothing about whether the model that combination lives in is the
right one. Coupling the reading to a misspecification diagnostic---so that one knows
both that a direction is identified and that the model along it is trustworthy---is
the natural next requirement for credible counterfactuals, and the harder one \est{open
question}.

\paragraph{Scaling}
The $O(P^3)$ eigendecomposition is negligible at $P=5$ and will not be at $P=10^3$
\est{established}. Whether the constrained subspace can be tracked without forming
$\Fhat$ in full---through matrix-free products with the gradient cloud, or a low-rank
sketch maintained along the path---is the engineering question standing between the
present validation and application to the large ABMs the differentiable programme
targets \est{open question}.

% =====================================================================
\appendix
% =====================================================================
\section{Curvature Preconditioning}
\label{app:preconditioning}

The same $\Fhat$ that diagnoses the calibration can precondition it. This is a
by-product rather than a contribution of the paper, and we report it to show that the
curvature the diagnostic reads is quantitatively as well as qualitatively useful.
Here the scale of $\Fhat$ matters, so unlike in the eigenstructure reading we do not
discard it (Section~\ref{sec:bg:opg}).

Because $\Fhat$ estimates the Gauss--Newton curvature $J^\top J$, the step
\begin{equation}
  \theta_{t+1} = \theta_t - \eta\,(\Fhat_t + \lambda I)^{-1}\,\bar g_t
  \label{eq:precond}
\end{equation}
is a Levenberg--Marquardt update, with the damping $\lambda$ interpolating between a
Gauss--Newton step ($\lambda\to0$) and gradient descent ($\lambda\to\infty$) and
regularising the sloppy directions the diagnostic flags as near-singular
\citep{transtrum2011geometry, martens2020new}. The mean gradient $\bar g_t$ and the
matrix $\Fhat_t$ are already available at each iterate, so the preconditioned step
reuses the diagnostic's own quantities.

\begin{figure*}[t]
  \centering
  \includegraphics[width=\textwidth]{images/figA_preconditioning.png}
  \caption{\textbf{Curvature preconditioning (Brock--Hommes).} Best-so-far $\mathrm{MMD}^2$
  (top) and distance to the ground-truth parameter $\|\theta_t-\theta^\star\|$
  (bottom) against iteration, for the OPG-Levenberg--Marquardt step of
  Eq.~\eqref{eq:precond} against Adam, SGD and L-BFGS, at three initialisation
  distances $d_0$. Shaded bands are inter-quartile ranges over seeds. OPG-LM is
  competitive with the strongest baseline on the easy and medium starts and, like all
  four methods, does not recover the parameter from the hard start---an
  initialisation far enough that the local curvature reading no longer points toward
  $\theta^\star$.}
  \label{fig:preconditioning}
\end{figure*}

Figure~\ref{fig:preconditioning} compares Eq.~\eqref{eq:precond} against Adam, SGD and
L-BFGS at three initialisation distances. On the easy and medium starts OPG-LM reaches
the best-so-far loss of the strongest baseline and tracks or leads on distance to
$\theta^\star$; on the hard start no method recovers the parameter, the OPG-LM step
included. We report the hard-start failure because it is consistent with the rest of
the paper rather than despite it: the preconditioner is built from a local curvature
estimate, and from an initialisation far enough that the local geometry no longer
points toward $\theta^\star$, no local method should be expected to succeed. The
result is a proof of concept---the diagnostic's matrix accelerates the fit it
diagnoses when the fit is within local reach---and not a claim to a state-of-the-art
optimiser.

% =====================================================================
\section{Technical Details}
\label{app:technical}

\paragraph{Gradient horizon}
Figure~\ref{fig:horizon} reports $\log_{10}\lambda_k$ against the gradient horizon $H$,
the number of final simulation steps retained in the backward pass. For Brock--Hommes
the whole spectrum grows in magnitude with $H$ but preserves its ordering, so the
stiff/sloppy hierarchy---and hence the diagnostic's qualitative reading---is robust to
truncation. For mean-field SIR the spectrum collapses to the floor for every $H < T$
and rises only at the full horizon $H = T$: the epidemic is an early transient, the
truncated gradient retains only the post-peak tail, and the identifiability
information lives in the growth phase that truncation discards. The SIR eigenstructure
of Figure~\ref{fig:sir_contrast} is therefore reported at full horizon only, and the
sensitivity to $H$ is itself a finding---a transient observable and a truncated
gradient are jointly incompatible with the reading, and practitioners who truncate for
memory should not expect it to hold.

\paragraph{Surrogate gradient}
Figure~\ref{fig:surrogate} compares the spectrum and eigenvectors under the
Gumbel-Sigmoid and straight-through surrogates. The sloppy directions agree to
$3$--$6^\circ$ between the two; the stiff directions rotate by about $28^\circ$ and the
eigenvalue magnitudes shift by roughly two orders. The bias thus concentrates in the
best-identified directions, as discussed in Sections~\ref{sec:exp:sir}
and~\ref{sec:discussion}. We tested two surrogates and do not extrapolate to others.

\paragraph{Falsification, SIR}
Figure~\ref{fig:sir_falsification} is the SIR counterpart of
Figure~\ref{fig:bh_falsification}. The sloppy lockdown combination is suppressed by
$20$--$30\times$ relative to the stiff direction under all three non-MMD
discrepancies. The stiff direction's response is asymmetric because full epidemic
suppression saturates the discrepancy, as described in Section~\ref{sec:exp:sir}.

\paragraph{Experimental configuration}
Both models use $P=5$ parameters and are calibrated to synthetic data generated at a
known $\theta^\star$. The MMD uses a characteristic kernel with bandwidth set by the
median heuristic. The gradient second moment is formed from $M \ge 2P$ per-seed
gradients per iterate; bootstrap intervals use $B=500$ resamples; the precision floor
is $\tau_t = 10^{-8}\lambda_1(t)$. Brock--Hommes is calibrated from three
regime-specific initialisations; mean-field SIR is calibrated at full gradient horizon
$H=T$.\vc\ \textcolor{red}{(Author: fill exact $M$, kernel, learning rates, seed
counts, hardware from the run logs before release.)}

\begin{figure}[t]
  \centering
  \includegraphics[width=\linewidth]{images/fig6_sir_falsification.png}
  \caption{\textbf{Falsification (mean-field SIR).} As
  Figure~\ref{fig:bh_falsification}, for the SIR model: the sloppy lockdown
  combination is suppressed by $20$--$30\times$ relative to the stiff direction under
  all three non-MMD discrepancies. The stiff direction's asymmetric response reflects
  epidemic-suppression saturation and is a feature of the dynamics, not of the
  diagnostic.}
  \label{fig:sir_falsification}
\end{figure}

\begin{figure}[t]
  \centering
  \includegraphics[width=\linewidth]{images/figB1_horizon.png}
  \caption{\textbf{Gradient-horizon sensitivity.} $\log_{10}\lambda_k$ against
  gradient horizon $H$ for Brock--Hommes (left) and mean-field SIR (right).
  Brock--Hommes grows in magnitude but keeps its hierarchy; SIR collapses to the floor
  for every $H<T$ and rises only at the full horizon, so its eigenstructure is
  meaningful only at $H=T$.}
  \label{fig:horizon}
\end{figure}

\begin{figure}[t]
  \centering
  \includegraphics[width=\linewidth]{images/figB2_surrogate.png}
  \caption{\textbf{Surrogate-gradient robustness (mean-field SIR).} Spectra under the
  Gumbel-Sigmoid and straight-through surrogates (left) and the per-eigenvector angle
  between their eigenvectors (right). The sloppy subspace is robust to the surrogate
  ($3$--$6^\circ$); the stiff directions rotate by about $28^\circ$ and the
  magnitudes shift by roughly two orders. The bias concentrates in the best-identified
  directions.}
  \label{fig:surrogate}
\end{figure}

\clearpage

% =====================================================================
% BIBLIOGRAPHY
% NOTE (author-facing): entries marked % <-- VERIFY were introduced during
% the rewrite and are NOT from the original .bib. Confirm each against its
% DOI before release. The \vc dagger in the text marks cost/quantitative
% claims to re-check against your run logs and the cited sources.
% =====================================================================
{\small
\begin{thebibliography}{99}
\setlength{\itemsep}{1pt}

\bibitem[Andelfinger(2023)]{andelfinger2023towards}
Andelfinger, P. (2023). Towards differentiable agent-based simulation.
\textit{ACM Trans. Model. Comput. Simul.}, 32(4), Article~23.

\bibitem[Arya et al.(2022)]{arya2022automatic}
Arya, G., Schauer, M., Sch\"afer, F., \& Rackauckas, C. (2022). Automatic
differentiation of programs with discrete randomness. \textit{NeurIPS} 35.

\bibitem[Axtell \& Farmer(2022)]{axtell2022agent} % <-- VERIFY (added in rewrite)
Axtell, R.~L., \& Farmer, J.~D. (2022). Agent-based modeling in economics and
finance: Past, present, and future. \textit{J. Econ. Lit.} (survey; verify exact
year/venue).

\bibitem[Bengio et al.(2013)]{bengio2013estimating}
Bengio, Y., L\'eonard, N., \& Courville, A. (2013). Estimating or propagating
gradients through stochastic neurons for conditional computation.
\textit{arXiv}:1308.3432.

\bibitem[Bonabeau(2002)]{bonabeau2002agent} % <-- VERIFY (added in rewrite)
Bonabeau, E. (2002). Agent-based modeling: Methods and techniques for simulating
human systems. \textit{PNAS}, 99(suppl.~3):7280--7287.

\bibitem[Botev et al.(2017)]{botev2017practical}
Botev, A., Ritter, H., \& Barber, D. (2017). Practical Gauss--Newton optimisation
for deep learning. \textit{ICML}, 557--565.

\bibitem[Brock \& Hommes(1998)]{brock1998heterogeneous}
Brock, W.~A., \& Hommes, C.~H. (1998). Heterogeneous beliefs and routes to chaos in
a simple asset pricing model. \textit{J. Econ. Dyn. Control}, 22(8--9):1235--1274.

\bibitem[Ch\'erief-Abdellatif \& Alquier(2020)]{cherief2020mmd}
Ch\'erief-Abdellatif, B.-E., \& Alquier, P. (2020). MMD-Bayes: Robust Bayesian
estimation via maximum mean discrepancy. \textit{AABI}, 1--21.

\bibitem[Chopra et al.(2023)]{chopra2023differentiable}
Chopra, A., Rodr\'iguez, A., Subramanian, J., Quera-Bofarull, A., Krishnamurthy, B.,
Prakash, B.~A., \& Raskar, R. (2023). Differentiable agent-based epidemiology.
\textit{AAMAS}, 1848--1857.

\bibitem[Cobelli \& DiStefano(1980)]{cobelli1980parameter}
Cobelli, C., \& DiStefano, J.~J. (1980). Parameter and structural identifiability
concepts and ambiguities. \textit{Am. J. Physiol.}, 239(1):R7--R24.

\bibitem[Cranmer et al.(2020)]{cranmer2020frontier}
Cranmer, K., Brehmer, J., \& Louppe, G. (2020). The frontier of simulation-based
inference. \textit{PNAS}, 117(48):30055--30062.

\bibitem[Dyer et al.(2024)]{dyer2024black}
Dyer, J., Cannon, P., Farmer, J.~D., \& Schmon, S.~M. (2024). Black-box Bayesian
inference for agent-based models. \textit{J. Econ. Dyn. Control}, 161:104827.

\bibitem[Eisenberg \& Hayashi(2014)]{eisenberg2014determining}
Eisenberg, M.~C., \& Hayashi, M.~A.~L. (2014). Determining identifiable parameter
combinations using subset profiling. \textit{Math. Biosci.}, 256:116--126.

\bibitem[Efron(1979)]{efron1979bootstrap}
Efron, B. (1979). Bootstrap methods: Another look at the jackknife.
\textit{Ann. Stat.}, 7(1):1--26.

\bibitem[Farmer \& Foley(2009)]{farmer2009economy} % <-- VERIFY (added in rewrite)
Farmer, J.~D., \& Foley, D. (2009). The economy needs agent-based modelling.
\textit{Nature}, 460(7256):685--686.

\bibitem[Grazzini et al.(2017)]{grazzini2017bayesian}
Grazzini, J., Richiardi, M.~G., \& Tsionas, M. (2017). Bayesian estimation of
agent-based models. \textit{J. Econ. Dyn. Control}, 77:26--47.

\bibitem[Gretton et al.(2012)]{gretton2012kernel}
Gretton, A., Borgwardt, K.~M., Rasch, M.~J., Sch\"olkopf, B., \& Smola, A. (2012).
A kernel two-sample test. \textit{JMLR}, 13:723--773.

\bibitem[Gutenkunst et al.(2007)]{gutenkunst2007universally}
Gutenkunst, R.~N., Waterfall, J.~J., Casey, F.~P., Brown, K.~S., Myers, C.~R., \&
Sethna, J.~P. (2007). Universally sloppy parameter sensitivities in systems biology
models. \textit{PLoS Comput. Biol.}, 3(10):e189.

\bibitem[Jang et al.(2017)]{jang2017categorical}
Jang, E., Gu, S., \& Poole, B. (2017). Categorical reparameterization with
Gumbel-softmax. \textit{ICLR}.

\bibitem[Knoblauch et al.(2022)]{knoblauch2022optimization}
Knoblauch, J., Jewson, J., \& Damoulas, T. (2022). An optimization-centric view on
Bayes' rule. \textit{JMLR}, 23(132):1--109.

\bibitem[Kunstner et al.(2019)]{kunstner2019limitations}
Kunstner, F., Hennig, P., \& Balles, L. (2019). Limitations of the empirical Fisher
approximation for natural gradient descent. \textit{NeurIPS} 32.

\bibitem[Li et al.(2018)]{li2018visualizing}
Li, H., Xu, Z., Taylor, G., Studer, C., \& Goldstein, T. (2018). Visualizing the
loss landscape of neural nets. \textit{NeurIPS} 31.

\bibitem[Machta et al.(2013)]{machta2013parameter}
Machta, B.~B., Chachra, R., Transtrum, M.~K., \& Sethna, J.~P. (2013). Parameter
space compression underlies emergent theories and predictive models.
\textit{Science}, 342(6158):604--607.

\bibitem[Maddison et al.(2017)]{maddison2017concrete}
Maddison, C.~J., Mnih, A., \& Teh, Y.~W. (2017). The Concrete distribution: A
continuous relaxation of discrete random variables. \textit{ICLR}.

\bibitem[Martens(2020)]{martens2020new}
Martens, J. (2020). New insights and perspectives on the natural gradient method.
\textit{JMLR}, 21(146):1--76.

\bibitem[Naumann-Woleske et al.(2024)]{naumann2024exploration}
Naumann-Woleske, K., Knicker, M.~S., Benzaquen, M., \& Bouchaud, J.-P. (2024).
Exploration of the parameter space in macroeconomic agent-based models.
\textit{Routledge Handbook of Complexity Economics}.

\bibitem[Platt(2020)]{platt2020comparison}
Platt, D. (2020). A comparison of economic agent-based model calibration methods.
\textit{J. Econ. Dyn. Control}, 113:103859.

\bibitem[Quera-Bofarull et al.(2023)]{querabofarull2023bayesian}
Quera-Bofarull, A., Chopra, A., Calinescu, A., Wooldridge, M., \& Dyer, J. (2023).
Bayesian calibration of differentiable agent-based models. \textit{ICLR Workshop on
AI for ABM}; \textit{arXiv}:2305.15340.

\bibitem[Quera-Bofarull et al.(2025)]{querabofarull2025automatic}
Quera-Bofarull, A., Bishop, N., Dyer, J., Jarne~Ornia, D., Calinescu, A., Farmer,
J.~D., \& Wooldridge, M. (2025). Automatic differentiation of agent-based models.
\textit{arXiv}:2509.03303.

\bibitem[Raue et al.(2009)]{raue2009structural}
Raue, A., Kreutz, C., Maiwald, T., Bachmann, J., Schilling, M., Klingm\"uller, U.,
\& Timmer, J. (2009). Structural and practical identifiability analysis of partially
observed dynamical models by exploiting the profile likelihood.
\textit{Bioinformatics}, 25(15):1923--1929.

\bibitem[Rothenberg(1971)]{rothenberg1971identification}
Rothenberg, T.~J. (1971). Identification in parametric models.
\textit{Econometrica}, 39(3):577--591.

\bibitem[Scarrold(2024)]{scarrold2024computing}
Scarrold, W. (2024). \textit{Computing Parameter Sloppiness in Nondeterministic
Economic Models}. DPhil thesis, University of Oxford.

\bibitem[Sriperumbudur et al.(2010)]{sriperumbudur2010hilbert}
Sriperumbudur, B.~K., Gretton, A., Fukumizu, K., Sch\"olkopf, B., \& Lanckriet,
G.~R.~G. (2010). Hilbert space embeddings and metrics on probability measures.
\textit{JMLR}, 11:1517--1561.

\bibitem[Transtrum et al.(2010)]{transtrum2010why}
Transtrum, M.~K., Machta, B.~B., \& Sethna, J.~P. (2010). Why are nonlinear fits to
data so challenging? \textit{Phys. Rev. Lett.}, 104(6):060201.

\bibitem[Transtrum et al.(2011)]{transtrum2011geometry}
Transtrum, M.~K., Machta, B.~B., \& Sethna, J.~P. (2011). Geometry of nonlinear
least squares with applications to sloppy models and optimization.
\textit{Phys. Rev. E}, 83(3):036701.

\bibitem[Wieland et al.(2021)]{wieland2021structural}
Wieland, F.-G., Hauber, A.~L., Rosenblatt, M., T\"onsing, C., \& Kreutz, C. (2021).
On structural and practical identifiability. \textit{Curr. Opin. Syst. Biol.},
25:60--69.

\end{thebibliography}
}

\end{document}